\clearpage
\subsection{Open Boundary Condition Case}

When $g_I=0$, we have $g_F=0$, and the fermionic model has open boundary conditions (OBC). In this case, translational invariance is broken, and we cannot apply the Fourier transformation used for periodic systems. Instead, we use the Bogoliubov-de Gennes (BdG) formalism and the analytic diagonalization approach developed by Lieb, Schultz, and Mattis.

Starting from the same Jordan--Wigner convention as in \cref{eq:H_ferm_realspace}, the open-chain fermionic Hamiltonian is
\begin{align}
    H_{\mathrm{spin}}^{\mathrm{OBC}}
        &= \frac{\cos\theta}{2}\sum_{n=1}^{N-1}
           (a_n-a_n^\dag)(a_{n+1}+a_{n+1}^\dag) \nn\\
        &\quad
           -\frac{\sin\theta}{2}\sum_{n=1}^{N}
           (a_n a_n^\dag-a_n^\dag a_n).
    \label{eq:H_obc_fermion_canonical}
\end{align}
Equivalently, after normal ordering the onsite term, the coefficient of $a_n^\dag a_n$ is $+\sin\theta$.  Thus, if $s=\sin\theta$ and $c=\cos\theta$, the BdG matrices are
\begin{align}
    A_{nn} &= s, &
    A_{n,n+1}=A_{n+1,n} &= -\frac{c}{2}, \nn\\
    B_{n,n+1} &= -\frac{c}{2}, &
    B_{n+1,n} &= \frac{c}{2},
    \label{eq:obc_AB_canonical}
\end{align}
with all other entries zero.  The matrix $A$ is symmetric and $B$ is antisymmetric.  In the Nambu spinor
\begin{equation}
    \bm{\Psi}=(a_1,\ldots,a_N,a_1^\dag,\ldots,a_N^\dag)^T
\end{equation}
the quadratic Hamiltonian is represented by
\begin{align}
    H_{\mathrm{BdG}} &=
    \begin{pmatrix}
        A  &B\\
        -B &-A
    \end{pmatrix}.
    \label{eq:obc_bdg_matrix}
\end{align}
The positive eigenvalues of this real symmetric matrix are the open-chain quasiparticle energies $\epsilon_k$.
To diagonalize $H_{\mathrm{spin}}^{\mathrm{OBC}}$, we use a Bogoliubov transformation to find new fermionic operators $\gamma_k, \gamma_k^\dag$ that are linear combinations of $a_n, a_n^\dag$:
\begin{align}
    \gamma_k      &= \sum_{n=1}^N (u_{nk} a_n + v_{nk} a_n^\dag)\\
    \gamma_k^\dag &= \sum_{n=1}^N (u_{nk}^* a_n^\dag + v_{nk}^* a_n)
\end{align}
With these operators $\gamma_k$ the Hamiltonian should become diagonal, and the eigenvalues should be the mode energies $\epsilon_k$:
\begin{equation}
    H = \sum_k \epsilon_k \left(\gamma_k^\dagger \gamma_k - \frac{1}{2}\right).
\end{equation}
which should give the commutation relations:
\begin{equation}
    [H, \gamma_k] = -\epsilon_k \gamma_k.
\end{equation}
such that we can derive the Bogoliubov-de Gennes eigenvalue equations:
\begin{align}
    A u_k + B v_k  &= \epsilon_k u_k \label{eq:BdG_eq1}\\
    -B u_k - A v_k &= \epsilon_k v_k \label{eq:BdG_eq2}
\end{align}
where $u_k = (u_{1k}, \dots, u_{Nk})^T$ and $v_k = (v_{1k}, \dots, v_{Nk})^T$ are $N$-component vectors containing the coefficients of the transformation.

Following Lieb, Schultz, and Mattis, we introduce new vectors $\phi_k = u_k + v_k$ and $\psi_k = u_k - v_k$. Adding and subtracting the BdG equations \cref{eq:BdG_eq1,eq:BdG_eq2}, we get:
\begin{align}
    (A - B) \psi_k &= \epsilon_k \phi_k \label{eq:LSM_eq1}\\
    (A + B) \phi_k &= \epsilon_k \psi_k \label{eq:LSM_eq2}
\end{align}
Eliminating $\psi_k$ (or $\phi_k$), we obtain:
\begin{align}
    (A - B)(A + B) \phi_k &= \epsilon_k^2 \phi_k \label{eq:E2_phi}\\
    (A + B)(A - B) \psi_k &= \epsilon_k^2 \psi_k \label{eq:E2_psi}
\end{align}
The key insight from the Lieb-Schultz-Mattis approach is that for our specific matrices $A$ and $B$, the products $(A+B)(A-B)$ and $(A-B)(A+B)$ have tridiagonal form. Calculating these products explicitly:
\begin{align}
    A+B &=
    \begin{pmatrix}
        s      & -c     &        &        \\
               & s      & \ddots &        \\
               &        & \ddots & -c     \\
               &        &        & s
    \end{pmatrix},
    &
    A-B &=
    \begin{pmatrix}
        s      &        &        &        \\
        -c     & s      &        &        \\
               & \ddots & \ddots &        \\
               &        & -c     & s
    \end{pmatrix}.
\end{align}
Consequently, the matrix acting on $\phi_k$ is
\begin{align}
    (A-B)(A+B) =
    \begin{pmatrix}
        s^2 & -sc &        &        \\
        -sc & 1   & \ddots &        \\
            & \ddots & \ddots & -sc \\
            &        & -sc & 1
    \end{pmatrix},
    \label{eq:obc_phi_tridiagonal}
\end{align}
while $(A+B)(A-B)$ has the same off-diagonal entries and diagonal entries $1,\ldots,1,s^2$.  The bulk eigenvalue equation for $\phi_{k,n}$ is therefore
\begin{align}
    \epsilon_k^2 \phi_{k,n}
        = \phi_{k,n}
          -sc(\phi_{k,n-1}+\phi_{k,n+1}),
\end{align}
and hence
\begin{align}
    \epsilon_k^2 = 1 - 2sc\cos\omega_k
                 = 1-\sin(2\theta)\cos\omega_k.
    \label{eq:obc_epsilon_canonical}
\end{align}
This tridiagonal structure allows for an analytic solution. The eigenvectors of these matrices have the form:
\begin{align}
    \phi_{k,n} &= C_{\phi,k}\sin((N+1-n)\omega_k),\\
    \psi_{k,n} &= C_{\psi,k}\sin(n\omega_k),
\end{align}
where \cref{eq:LSM_eq2} fixes the relative normalization, for example
\begin{equation}
    \frac{C_{\psi,k}}{C_{\phi,k}}
        = \frac{s\sin\omega_k}{\epsilon_k\sin(N\omega_k)}
\end{equation}
away from the limiting edge-mode cases, and the remaining overall normalization is fixed by the canonical anticommutation relation for $\gamma_k$.  The first row of \cref{eq:obc_phi_tridiagonal} imposes the open-boundary quantization condition
\begin{align}
    \frac{\sin((N+1)\omega_k)}{\sin(N\omega_k)} = \cot\theta,
    \label{eq:obc_quantization_canonical}
\end{align}
with the usual limiting treatment when a root is exponentially close to an edge-mode value.  Combining \cref{eq:obc_epsilon_canonical,eq:obc_quantization_canonical} gives the mode energies
\begin{align}
    E_k = \epsilon_k = \sqrt{1-\sin(2\theta)\cos\omega_k}.
\end{align}
The diagonalized Hamiltonian takes the form:
\begin{align}
    H_{\mathrm{spin}}^{\mathrm{OBC}} &= \sum_{k=1}^N \epsilon_k \left(\gamma_k^\dag \gamma_k - \frac{1}{2}\right)
\end{align}

\subsection{The numerical method}

There is a more direct numerical method, and this is the convention used in the code.  One forms the matrices $A$ and $B$ in \cref{eq:obc_AB_canonical} and diagonalizes the real symmetric BdG matrix \cref{eq:obc_bdg_matrix}.  Its eigenvalues occur in pairs $-\epsilon_k,\epsilon_k$, and the positive half gives the quasiparticle spectrum.  In the Julia implementation this source of truth is \texttt{obc\_bdg\_matrices} and \texttt{obc\_mode\_energies} in \texttt{src/mode\_analysis.jl}.

If $W$ is an orthogonal matrix of BdG eigenvectors, then the transformed Nambu vector gives the Bogoliubov operators $b_k,b_k^\dagger$.  The Hamiltonian is diagonal in this basis:
\begin{equation}
    H = \sum_k \epsilon_k \left(b_k^\dagger b_k - \frac{1}{2}\right).
\end{equation}
The many-body spectrum is therefore obtained from the subset sums $\sum_k \epsilon_k(n_k-\frac12)$ with $n_k\in\{0,1\}$.  The regression tests compare this spectrum, after multiplying by the code scale $\Lambda=2\sqrt{J^2+h^2}$, with the exact diagonalization spectrum of the open spin Hamiltonian.

\subsection{Mode energy and occupation}

To express the mode occupation operators $\gamma_k^\dag \gamma_k$ in terms of the original fermionic operators, we use the Bogoliubov coefficients derived from the eigenvectors:
\begin{align}
    u_{nk} &= \frac{1}{2}\left[
        C_{\phi,k}\sin((N+1-n)\omega_k)+C_{\psi,k}\sin(n\omega_k)\right],\\
    v_{nk} &= \frac{1}{2}\left[
        C_{\phi,k}\sin((N+1-n)\omega_k)-C_{\psi,k}\sin(n\omega_k)\right].
\end{align}
The mode occupation operator is
\begin{align}
    \hat n_k^{\mathrm{OBC}}
    \equiv \gamma_k^\dag \gamma_k
    &= \sum_{n,m=1}^N \Bigl[
        u_{nk}^* u_{mk} a_n^\dag a_m
        + u_{nk}^* v_{mk} a_n^\dag a_m^\dag \nn\\
    &\quad
        + v_{nk}^* u_{mk} a_n a_m
        + v_{nk}^* v_{mk} a_n a_m^\dag
    \Bigr].
\end{align}
Using the anticommutation relation $a_n a_m^\dag = \delta_{nm} - a_m^\dag a_n$, and applying the JW transformation from \cref{eq:JW_a_n,eq:JW_a_n_dag}, we can express $\gamma_k^\dag \gamma_k$ in terms of spin operators.  We use the normalized Pauli ladder operators
\begin{equation}
    \sigma_n^\pm=\frac{\sigma_{n,x}\pm i\sigma_{n,y}}{2}
\end{equation}
and keep the Jordan--Wigner strings in the split order in which the fermionic bilinears generate them:
\begin{widetext}
    \begin{align}
        \hat n_k^{\mathrm{OBC}}
        &= \sum_{n,m=1}^N \left[
            u_{nk}^* u_{mk}\, S_n\sigma_n^+ S_m\sigma_m^-
            + u_{nk}^* v_{mk}\, S_n\sigma_n^+ S_m\sigma_m^+
        \right. \nn\\
        &\quad \left.
            + v_{nk}^* u_{mk}\, S_n\sigma_n^- S_m\sigma_m^-
            - v_{nk}^* v_{mk}\, S_m\sigma_m^+ S_n\sigma_n^-
        \right]
        + \sum_{n=1}^N |v_{nk}|^2 .
        \label{eq:obc_mode_occupation_spin}
    \end{align}
\end{widetext}
Equivalently, the product of strings in each split-string bilinear reduces to
\begin{align}
    S_n S_m =
    \begin{cases}
        \prod_{j=\min(n,m)}^{\max(n,m)-1} \sigma_{j,z} & \text{if } n \neq m \\
        1                                              & \text{if } n = m
    \end{cases}
\end{align}
provided the operator ordering in \cref{eq:obc_mode_occupation_spin} is retained.

The dimensionless open-chain mode observable is
\begin{equation}
    h_k^{\mathrm{OBC}} = 2\hat n_k^{\mathrm{OBC}}-1 .
\end{equation}
With the diagonal convention above, the Hamiltonian contribution of this mode is
\begin{equation}
    H_k
    = \epsilon_k\left(\hat n_k^{\mathrm{OBC}}-\frac{1}{2}\right)
    = \frac{\epsilon_k}{2} h_k^{\mathrm{OBC}} .
    \label{eq:obc_mode_contribution}
\end{equation}
Thus \(H_k=\frac{\epsilon_k}{2}h_k^{\mathrm{OBC}}\), not \(\epsilon_k h_k^{\mathrm{OBC}}\).  This is the open-chain analogue of the factor \(1/2\) in the periodic full-grid reconstruction.
